\begin{abstract}
This study examined whether survey questions on generative artificial intelligence (GenAI) use and dependence improved estimates of current student burnout beyond study conditions.
The analysis included 97 cross-sectional survey responses, 87 from students who reported attending universities in Lebanon.
Burnout was measured on a 1-to-5 scale using the 12-item Burnout Assessment Tool for Students.
Both prediction tasks were evaluated using repeated nested cross-validation.
In matched comparisons using the same model type and test students, adding three GenAI questions to five study-condition questions reduced mean absolute error from 0.6183 to 0.5883 points.
Adding all GenAI questions to a broader set of background and study-condition questions reduced error from 0.6145 to 0.5933 points.
A separate comparison evaluated two approaches to identifying students with measured burnout scores of at least 4.
At the selected warning thresholds, the workload-only model detected an average of 16.0 of the 17 high-symptom cases, and the eight-question model detected 14.8.
Both averaged 28 false warnings among the 80 students below the cutoff. The conditional 95\% intervals for paired performance differences included zero.
The eight-question model achieved 76.0\% balanced accuracy, exceeding all 199 shuffled-label results (permutation \(p=0.005\)).
Independent validation is needed to assess whether the findings generalize to other student cohorts.
\end{abstract}
